\section{Joint hereditary selectors by finite descending filters}
\label{sec:typed-selectors}

We continue with the original (\emph{not} a contracted image)
strong exact families $\mathcal H_{ij}$ from
Section~\ref{sec:typed-strong}.  At no point do we declare a new
product on an intermediate image of a selector.  The hypotheses of
this section are precisely finite named symmetric typed palettes,
SAP, semantic homomorphic CCO, and
\begin{equation}\label{eq:NR2}
 (\mathrm{NR2})\qquad
 \text{no named sort admits a proper relative-index-two pair of}
 \text{ invariant equivalences}.
\end{equation}
The primitive named sorts of $U$ satisfy (NR2) by
Theorem~\ref{thm:joint-transversal}; nevertheless we spell out its
\emph{only} use in the proof.

\begin{definition}[Joint pruning maps]
\label{def:joint-pruning}
A \emph{joint hereditary pruning} $R$ assigns to every exact
$D\in\mathcal H_{ij}$ a nonempty (possibly nonexact) subset
$R_{ij}(D)\subseteq D$.  It is \emph{converse-compatible} if
$R_{ji}(\bar D)=\overline{R_{ij}(D)}$; it \emph{preserves PC} if
all three inclusions in \eqref{eq:typed-PC} hold for its three
outputs whenever they hold for the input; and it has \emph{survivor
heredity} if
\begin{equation}\label{eq:typed-survivor}
 E\subseteq D,\quad E\cap R(D)\ne\varnothing
 \quad\Longrightarrow\quad R(E)=E\cap R(D)
\end{equation}
whenever $D,E\in\mathcal H_{ij}$.  An \emph{exact contraction} is
such a map with each output in $\mathcal H_{ij}$.  A \emph{joint
hereditary selector} (HS) is a singleton-valued joint hereditary
pruning; write $R(D)=\{h_{ij}(D)\}$.

The implication \eqref{eq:typed-survivor} is an additional axiom:
conservativity, converse and triangle legality alone do not imply
it.  In the proof below every intermediate map is an exact
contraction and hence its output belongs to the \emph{original}
$\mathcal H$, while the final selector chooses atomic elements.
\end{definition}

\begin{definition}[Original multiplication-closed filters]
\label{def:typed-filter}
A \emph{current family} $\mathcal J=(\mathcal J_{ij})$ is a
converse-closed subfamily of $\mathcal H$ containing all singleton
atomic menus and downward closed \emph{relative to the original exact
family}: if $D\in\mathcal J_{ij}$ and
$E\in\mathcal H_{ij}$ with $\varnothing\ne E\subseteq D$,
then $E\in\mathcal J_{ij}$.  A \emph{filter family}
$\mathfrak F=(\mathfrak F_{ij})$ is a (possibly empty) tagged
subfamily of the \emph{original} exact menus, closed under converse
and all compatible left and right products by members of
$\mathcal H$.  Its \emph{relative visibility} on $\mathcal J$ is
\begin{equation}\label{eq:relative-visibility}
 \begin{gathered}
 D\in\mathcal J_{ij},\quad A,B\in\mathfrak F_{ij},
 \quad A\cap D\ne\varnothing,\ B\cap D\ne\varnothing\\
 \Longrightarrow A\cap B\ne\varnothing.
 \end{gathered}
\end{equation}
Given such a filter, define the gate on $D\in\mathcal J_{ij}$ by
\begin{equation}\label{eq:typed-gate}
 \mathcal G(D)=D\cap\bigcap\{A\in\mathfrak F_{ij}:
                                    A\cap D\ne\varnothing\}.
\end{equation}
The empty active list leaves $D$ unchanged; no intersections of
menus from different hom-palettes are formed.
\end{definition}

\begin{lemma}[The bilateral-support implication]
\label{lem:typed-bilateral}
Suppose $\mathfrak F$ is multiplication/converse closed and has
relative visibility on $\mathcal J$.  For
$D\in\mathcal J_{ij}$, $E\in\mathcal J_{kj}$ and a fixed
$P\in\mathcal H_{ik}$, assume the two \emph{original} inclusions
\begin{equation}\label{eq:typed-bilateral-assumption}
 D\subseteq P*E,\qquad E\subseteq\bar P*D.
\end{equation}
Then the gate from \eqref{eq:typed-gate} preserves both:
\begin{equation}\label{eq:typed-bilateral-conclusion}
 \mathcal G(D)\subseteq P*\mathcal G(E),\qquad
 \mathcal G(E)\subseteq\bar P*\mathcal G(D).
\end{equation}
The symmetric assertion for a fixed \emph{right} factor follows by
converse.  There is no assumption that $P$ belongs to the current
family $\mathcal J$.
\end{lemma}
\begin{proof}
Consider any filter $A\in\mathfrak F_{kj}$ active on $E$.
Choose $e\in A\cap E$.  The second inclusion in
\eqref{eq:typed-bilateral-assumption} gives a supported
$p\in P,d\in D$ with $e\in\bar p*d$, equivalently
$d\in p*e$.  Therefore $P*A\in\mathfrak F_{ij}$ is active
on $D$.  By definition of the gate,
$\mathcal G(D)\subseteq P*A$ for every such active $A$, and also
$\mathcal G(D)\subseteq D\subseteq P*E$.  Distributivity
\eqref{eq:typed-distrib} yields
\[
 \mathcal G(D)\subseteq
 (P*E)\cap\bigcap_{A\text{ active on }E}(P*A)
 =P*\biggl(E\cap\bigcap_{A\text{ active on }E}A\biggr)
 =P*\mathcal G(E).
\]
This includes an empty active list, with the evident top convention.
Interchanging $D,E$ and replacing $P$ with $\bar P$ proves the
second inclusion.  The argument with a fixed right factor is its
converse.  Every product used here is the original strong exact
product from Section~\ref{sec:typed-strong}.
\end{proof}

\begin{lemma}[All three PC coordinates survive one gate]
\label{lem:typed-all-PC}
Let $\mathcal J,\mathfrak F$ satisfy the hypotheses of
Definition~\ref{def:typed-filter} and
\eqref{eq:relative-visibility}.  Then the gate $\mathcal G$
of \eqref{eq:typed-gate} has nonempty exact outputs, commutes with
converse, satisfies survivor heredity and preserves
\emph{all three} support inclusions in every typed PC triple.
The statement includes $i=j$, $j=k$, $i=k$ and coinciding source
menus; no injectivity assumption on the variables denoting menus
is made.
\end{lemma}
\begin{proof}
For a fixed $D$ the active filters together with $D$ form a
pairwise-intersecting list by visibility.  Exact $2$-Helly
(Lemma~\ref{lem:typed-helly}) makes their intersection nonempty;
exact-meet closure places $\mathcal G(D)$ in $\mathcal H$ and
$\mathcal J$.  Conversely oriented active lists are exchanged by
bar, so $\mathcal G(\bar D)=\overline{\mathcal G(D)}$.

If $E\subseteq D$ and $E\cap\mathcal G(D)\ne\varnothing$, every
filter active on $D$ meets $E$; conversely every filter active on
$E$ is active on $D$.  Their active lists coincide.  Hence
\[
 \mathcal G(E)=E\cap\mathcal G(D),
\]
proving survivor heredity, including the equality and singleton
cases.

It remains to verify the \emph{first} PC inclusion; the other two
then follow by applying the same proof to the two converse/Peirce
rotations \eqref{eq:typed-PC-rotation}.
Let $\mathrm{PC}_{ijk}(A,B,C)$ and fix an \emph{arbitrary}
$c\in\mathcal G(C)$.  The fixed bilateral and right-support
implications of Lemma~\ref{lem:typed-bilateral}, applied to the
original PC inclusions, give
\[
 c\in A*\mathcal G(B),\qquad c\in\mathcal G(A)*B.
\]
In the \emph{original} exact families set
\[
 A_c=A\cap(\{c\}*\bar B),\qquad
 B_c=B\cap(\bar A*\{c\}).
\]
Both are nonempty exact submenus: the displayed memberships give
respectively an atom in $A_c\cap\mathcal G(A)$ and in
$B_c\cap\mathcal G(B)$.  They belong to $\mathcal J$ by its
downward closure, and survivor heredity gives
\begin{equation}\label{eq:typed-PC-restrictions}
 \mathcal G(A_c)=A_c\cap\mathcal G(A),\qquad
 \mathcal G(B_c)=B_c\cap\mathcal G(B).
\end{equation}
We now check \emph{both} sides of a fresh bilateral-support
condition with the singleton \emph{fixed multiplier} $\{c\}$.
By exact distributivity and the duplication inclusion
$X\subseteq\{c\}*\bar{\{c\}}*X$ (proved by using SAP for a
single supported fork and then applying Peirce),
\[
 \{c\}*\overline{B_c}
  =(\{c\}*\bar B)\cap
    (\{c\}*\bar{\{c\}}*A)
 \supseteq A_c.
\]
For the other direction, apply the same distributivity and
duplication to $\overline{B_c}=\bar B\cap
(\bar{\{c\}}*A)$:
\[
 \bar{\{c\}}*A_c
  =(\bar{\{c\}}*A)\cap
    (\bar{\{c\}}*\{c\}*\bar B)
 \supseteq\overline{B_c}.
\]
Thus $A_c\subseteq\{c\}*\overline{B_c}$ and
$\overline{B_c}\subseteq\bar{\{c\}}*A_c$.
Lemma~\ref{lem:typed-bilateral} applied to the pair
$(A_c,\overline{B_c})$ gives
\[
 \mathcal G(A_c)\subseteq\{c\}*\mathcal G(\overline{B_c}).
\]
Choose any $a\in\mathcal G(A_c)$, which is nonempty.  The last
inclusion gives $b\in\mathcal G(B_c)$ with
$c\in a*b$ by Peirce.  The containment
\eqref{eq:typed-PC-restrictions} implies
$a\in\mathcal G(A)$ and $b\in\mathcal G(B)$.
Consequently $c\in\mathcal G(A)*\mathcal G(B)$.
Since $c$ was arbitrary,
\[
 \mathcal G(C)\subseteq\mathcal G(A)*\mathcal G(B).
\]
The two rotated PC triples establish the inclusions for
$\mathcal G(A)$ and $\mathcal G(B)$.  The argument never assumes
that the menus or any of their endpoint sorts are different.
\end{proof}

\begin{lemma}[Image, composition and strict contraction]
\label{lem:typed-gate-image}
Under the preceding hypotheses, $\mathcal G$ fixes its image
pointwise; that image is a converse-closed downward subfamily of the
\emph{original} $\mathcal H$ and contains all atomic singletons.
Successive gates compose to exact conservative converse-compatible
PC-preserving maps with survivor heredity.  If some
$A\in\mathfrak F_{ij}$ has
$\varnothing\ne A\cap D\subsetneq D$ for a current menu
$D\in\mathcal J_{ij}$, then $\mathcal G(D)\subsetneq D$.
\end{lemma}
\begin{proof}
For any original exact $E\subseteq\mathcal G(D)$, survivor
heredity gives $\mathcal G(E)=E$.  In particular
$\mathcal G(\mathcal G(D))=\mathcal G(D)$; the image is exactly
the fixed family and is downward.  Atomic singletons remain fixed by
conservativity and nonemptiness.  Converse compatibility was proved
above.

For consecutive gates $G_1,G_2$, let $E\subseteq D$ meet
$G_2(G_1(D))$.  Then $G_1(E)=E\cap G_1(D)$ and this menu meets
$G_2(G_1(D))$.  Heredity of $G_2$ yields
\[
 G_2(G_1(E))=G_1(E)\cap G_2(G_1(D))
               =E\cap G_2(G_1(D)).
\]
Converse compatibility and PC preservation compose as well.
Finally an active proper trace occurs among the intersected filters
in \eqref{eq:typed-gate}, so it strictly reduces $D$.
\end{proof}

\begin{lemma}[Self-converse loop stabilization from (NR2)]
\label{lem:typed-loop-stable}
For every nonempty loop exact menu $P\in\mathcal H_{ii}$,
there exists $N$ such that $P^{*n}=P^{*(n+1)}$ for all $n\ge N$.
This is not a claim that the stable menu is the entire palette
$L_{ii}$.
\end{lemma}
\begin{proof}
All loop subsets are self-converse, so the product of two loop
menus commutes:
$P*Q=\overline{P*Q}=\bar Q*\bar P=Q*P$.
For any nonempty loop subset $P$, SAP and Peirce give the
\emph{duplication} $P\subseteq P^{*3}$: choose an auxiliary
$p\in P$, complete a fork witnessing $x\in p*(\bar p*x)$
for each $x\in P$.  Thus the even powers
$P^{*2}\subseteq P^{*4}\subseteq\cdots$ eventually stabilize
in the finite palette $L_{ii}$.  Write the stable even menu as $J$
and the stable odd menu as $O$.  Associativity and eventual
stability give
\begin{equation}\label{eq:even-odd}
 J*J=J,\quad J*O=O,\quad O*O=J,\quad P\subseteq O.
\end{equation}
For every $t\in J$ we have $J*\{t\}=J$:
the inclusion $\subseteq$ is closure under $J*J=J$;
for $r\in J$ take $u\in r*t$ by SAP and rotate the triple to
obtain $r\in u*t$ with $u\in J$.  If $J\cap O$ contains $t$,
then
$J=J*t\subseteq J*O=O$, and
$O=J*O\subseteq O*O=J$,
so $J=O$.  Otherwise the two nonempty disjoint loop menus
$J,O$ yield invariant equivalences
\[
 R=\Delta_i\cup J\quad\text{and}\quad
 S=\Delta_i\cup J\cup O
\]
on the named sort $i$, closed under triangle composition by
\eqref{eq:even-odd}.  For each point $x$, the $O$-neighbours
form precisely one other $R$-class: any two such neighbours
are $J$-related by $O*O=J$, and an $O$-neighbour exists by
realisation/transitivity.  Hence $[S:R]=2$, contradicting (NR2).
Therefore $J=O$ and all sufficiently high consecutive powers of
$P$ are equal.
\end{proof}

\begin{lemma}[Parallel-menu separation]
\label{lem:typed-no-return}
For any exact seed $C\in\mathcal H_{ab}$ let
$\mathfrak F(C)$ be the smallest \emph{tagged} family of original
exact menus containing $C$, closed under converse and all
compatible left/right multiplication by \emph{original} exact
menus.  If $A,B\in\mathcal H_{ij}$ are parallel menus and
$\mathfrak F(A)=\mathfrak F(B)$, then $A=B$.
\end{lemma}
\begin{proof}
Every menu generated from $C$ can be written in one of the two
normal forms
$P*C*Q$ or $P*\bar C*Q$, with outside multipliers
possibly \emph{absent}.  This follows by induction from the
closure rule and associativity; the two orientations are exchanged
by converse.  An absent multiplier means the identity \emph{action}
on menus, not an equality atom or an empty-word menu.

Suppose a generated menu $B\in\mathfrak F(A)$ has the \emph{same
ordered} endpoints $(i,j)$ as $A$.  A representation
$P*A*Q$ then has loop multipliers $P\in\mathcal H_{ii}$,
$Q\in\mathcal H_{jj}$, or absent actions.  If $i=j$, every loop
menu is self-converse and the reverse-oriented representation also
has this form.  If $i\ne j$, a representation
$P*\bar A*Q$ returning to $(i,j)$ has both
$P,Q\in\mathcal H_{ij}$, and neither is absent.  The loop
palettes are self-converse, so
\begin{equation}\label{eq:parallel-switch}
 P*\bar A=\overline{P*\bar A}=A*\bar P,
\end{equation}
and hence $P*\bar A*Q=A*(\bar P*Q)$; the latter right
multiplier is a loop menu of sort $j$.  Thus every oriented return
has the form $L*A*R$ with loop or absent multipliers.

Equal generated families supply two such returns
$T(A)=B$ and $V(B)=A$, where
$T(X)=L*X*R$ and $V(X)=P*X*Q$ on $\mathcal H_{ij}$.
They commute, because all relevant loop multipliers commute
inside their own endpoint palettes; no cross-arrow permutation is
used.  By Lemma~\ref{lem:typed-loop-stable} there exists $n>0$
with $T^n=T^{n+1}$ as maps, using an identity action for any
absent multiplier.  Since $VT(A)=A$, commutativity gives
$V^nT^n(A)=A$.  Therefore
\[
 B=T(A)=T\,V^nT^n(A)
       =V^nT^{n+1}(A)=V^nT^n(A)=A.
\]
This proves the assertion.  In particular the proof never uses the
ill-typed product of two cross-sort arrows $C*C$.
\end{proof}

\begin{lemma}[Gate iteration from pure-tag no-return]
\label{lem:typed-pure-action-gate-engine}
Let $\Lambda$ be a finite group permuting the named sorts of the
original strong exact typed system $\mathcal H$.  Assume this is a
\emph{pure tag action}: it fixes literal original fine-colour names
and transports all legal atomic triangles, converse, exact menus
and typed strong multiplication to the relabelled sort signatures.
For an original exact seed $C$, let $\mathfrak F_\Lambda(C)$
be the smallest family of \emph{original tagged exact menus}
containing $C$ and closed under converse, compatible left and
right multiplication by arbitrary original exact menus, and the
$\Lambda$ action.  Suppose that for every seed $C$ and every
parallel $A,B\in\mathfrak F_\Lambda(C)$,
\[
 \mathfrak F_\Lambda(A)=\mathfrak F_\Lambda(B)
      \quad\Longrightarrow\quad A=B
\]
as literal subsets of their original fine-colour palette.
Then the whole original $\mathcal H$ admits one joint
singleton-valued hereditary selector preserving converse and
all three PC inclusions, and satisfying
\[
 h_{\lambda i,\lambda j}(D)=h_{ij}(D)
 \quad(\lambda\in\Lambda,\ D\in\mathcal H_{ij}),
\]
where the same literal menu $D$ is regarded in the relabelled
palette by the pure action.  The conclusion does not assert
covariance under an action which permutes actual fine-colour names.
\end{lemma}
\begin{proof}
Start with the original full current family $\mathcal J=\mathcal H$,
which is downward, converse-closed, $\Lambda$-invariant and
contains every tagged atomic singleton.  If a nonsingleton menu
remains, an atomic singleton $C\subsetneq D$ is a seed with a
nonempty proper trace on the current menu $D$.  Among all
\emph{original} exact seeds having a nonempty proper current trace,
choose $C$ minimizing the finite number $|\mathfrak F_\Lambda(C)|$.

We check relative visibility.  If it fails, there are a current
$D'$ and parallel $A,B\in\mathfrak F_\Lambda(C)$ with
$A\cap D'\ne\varnothing$, $B\cap D'\ne\varnothing$ but
$A\cap B=\varnothing$.  Then each of $A,B$ is itself a seed
with a nonempty proper trace on $D'$, since the other one meets
$D'$ disjointly.  Closure under all original generators implies
$\mathfrak F_\Lambda(A),\mathfrak F_\Lambda(B)
\subseteq\mathfrak F_\Lambda(C)$.  If both inclusions were
equalities, the no-return hypothesis would force $A=B$, a
contradiction.  Thus one of the two viable seeds generates a
\emph{strictly} smaller family, contrary to minimality.

Apply the original gate \eqref{eq:typed-gate} to
$\mathfrak F_\Lambda(C)$.  By
\cref{lem:typed-all-PC,lem:typed-gate-image}, its output is a
nonempty exact, downward current family containing every tagged
singleton, and it preserves converse, survivor heredity and all
three typed PC coordinates.  A proper trace forces strict
contraction.  Because the generated family and current family
are $\Lambda$-invariant and $\Lambda$ fixes actual atom names,
the active filter lists are carried bijectively to those for the
relabelled menu; hence the gate itself is $\Lambda$-covariant.
Its image is therefore again a $\Lambda$-invariant current family.

Iterate until no nonsingleton menu remains.  The number of tagged
original exact menus in the current family decreases strictly,
so after at most
$\sum_{i,j}|\mathcal H_{ij}|-\sum_{i,j}|L_{ij}|$
proper gates the process terminates.  Composition retains
converse, PC, survivor heredity and $\Lambda$ covariance by
\cref{lem:typed-gate-image}.  Since the terminal exact menus are
atomic singletons, the composite map is the asserted selector.
The choice of seeds is finite and need not be $\Lambda$-invariant:
each whole $\Lambda$-closed generated family is used by its gate.
\end{proof}

\begin{theorem}[Existence of a joint hereditary selector]
\label{thm:typed-HS}
Let $\mathcal K$ be any finite-sorted symmetric complete atomic
binary triangle Fra\"iss\'e class with SAP, semantic homomorphic
CCO, and (NR2) on every named sort.  Its full strong exact family
$\mathcal H=(\mathcal H_{ij})$ admits a joint hereditary selector
$h$: it chooses $h_{ij}(D)\in D$, commutes with converse, preserves
all typed PC triples, and satisfies
\begin{equation}\label{eq:HS-heredity}
 E\subseteq D,\ h_{ij}(D)\in E\quad\Longrightarrow\quad
 h_{ij}(E)=h_{ij}(D).
\end{equation}
No common identification of the palettes $L_{ij}$ is required.
In particular this theorem applies to the actual joint primitive
transversal from Theorem~\ref{thm:joint-transversal}.
\end{theorem}
\begin{proof}
Apply the general gate-iteration lemma
\cref{lem:typed-pure-action-gate-engine} with the \emph{trivial}
tag group $\Lambda=\{1\}$.  Its generated family
$\mathfrak F_\Lambda(C)$ is exactly the original
$\mathfrak F(C)$ generated family used in
Lemma~\ref{lem:typed-no-return}; that same lemma
supplies its required no-return implication.  Thus all hypotheses
of the common gate mechanism hold under the present strong
exact/NR2 assumptions, and its singleton output has the
required PC, converse and survivor-heredity properties.
No covariance under any nontrivial type permutation is claimed.
\end{proof}

\begin{remark}[Scope of the selector theorem]
\label{rem:typed-selector-scope}
The theorem establishes \emph{one joint selector} for the fixed
\emph{named} typed source and supports all three PC coordinates,
including parallel or repeated menus and sort labels.  It does
\emph{not} imply that the selector is covariant under an arbitrary
permutation of owner names, even if such a permutation preserves all
atomic triangle tables.  The additional pure-owner action and its
covariant selector descent are treated separately in the next
construction chapter; they cannot be obtained by naively averaging
selectors over a group.
\end{remark}
